\documentclass[legal]{polinetwork}

\pnTitle{Privacy Policy}
\pnSubtitle{Administrators and Volunteers}

\begin{document}

\pnMakeHeader

\section{Data Controller}

Pursuant to Articles 4 and 24 of EU Regulation 2016/679, the Data Controller is the legal representative of the association PoliNetwork APS, i.e., its President.

The association has not appointed a Data Protection Officer (DPO) as the conditions set out under Article 37 of the GDPR are not met.

\textbf{Contact details to exercise privacy rights:} \\

\textbf{Email address:} \href{mailto:privacy@polinetwork.org}{privacy@polinetwork.org} (alternatively \href{mailto:direttivo@polinetwork.org}{direttivo@polinetwork.org}) \\
\textbf{PEC (Posta Elettronica Certificata – Italian Certified Electronic Mail):} \href{mailto:polinetwork@pec.it}{polinetwork@pec.it}

\section{Types of Data Collected}

The Data Controller processes the personal identification data communicated by the data subject when applying for a volunteer position at PoliNetwork APS:

\begin{itemize}
	\item First name, Last name, Any alias;
	\item Email address, Phone number;
	\item Username set on the Telegram platform;
	\item Start date of the volunteering activity;
	\item Information regarding studies, any representative role at the Politecnico di Milano, and roles within the association and/or other associations;
	\item Any additional information provided by the volunteer;
	\item Tax Identification Number (Codice Fiscale) (required for members of the HR team).
\end{itemize}

\section{Purposes and Legal Basis for Processing}

The personal data provided will be processed in compliance with the lawfulness conditions set out in EU Regulation 2016/679 (in particular pursuant to Art. 6(b), (c) and (f)), for the establishment, organisational management, and planning of the collaborative relationship as a volunteer for the association, and in particular for:

\begin{itemize}
	\item entry into internal records, assignment of operational roles, and management of related permissions or access;
	\item coordination of activities, task management, and internal communication necessary for carrying out the role;
	\item internal organisational archiving and analysis, including monitoring the composition of working groups and assessing the availability of collaborators (for example, in anticipation of the completion of their studies), in order to plan and launch new recruitment campaigns promptly;
	\item fulfilling contractual and legal obligations and administrative-accounting purposes. For the purposes of applying data protection legislation, processing carried out for administrative-accounting purposes refers to those connected with the organisation, administration, and management of collaborators;
	\item fulfilling obligations required by law (such as taking out mandatory insurance coverage for volunteers and maintaining the registers of the Third Sector), by a regulation, by EU legislation, or by an order from the Authority;
	\item exercising the rights of the Data Controller, for example the right of defence in legal proceedings.
\end{itemize}

The provision of personal data for the purposes listed above is necessary to proceed with establishing the collaborative relationship. Failure to provide the required personal data may result in the inability to assume the role of administrator or volunteer within the association.

\section{Methods of Processing}

The processing of personal data is carried out using manual, computerised, and telematic tools. The logic applied is strictly related to the purposes themselves and implemented in such a way as to ensure the maximum security and confidentiality of the data, preventing loss, unlawful use, and unauthorised access.

\section{Data Retention}

Processing will be carried out in automated and/or manual form, using methods and tools designed to ensure maximum security and confidentiality, by persons specifically appointed for this purpose. In compliance with Article 5(1)(e) of EU Regulation 2016/679, the personal data collected will be stored in a form that allows the identification of data subjects for no longer than is necessary for the purposes for which the personal data are processed.

\section{Recipients and Access to Data}

Personal data will not be communicated to external third parties or publicly disclosed. Data will be accessible exclusively to members of the PoliNetwork APS association directly involved in the administrative, organisational, and operational management of volunteers (for example, members of the Board of Directors, HR area contacts, or project coordinators with whom the volunteer is involved). These persons have been specifically instructed and will process the data in full compliance with internal confidentiality guidelines.

Access to the data may also be granted to those who provide services for the management and maintenance of the Data Controller’s IT systems, who act as Data Processors pursuant to Article 28 of the GDPR. In particular, Microsoft Forms, provided by Microsoft Corporation, is used to collect volunteer data and registrations.

\textbf{Transfer of Data to Third Countries.} Microsoft Corporation may process personal data in the United States of America. Such transfers take place in compliance with the GDPR on the basis of the following safeguards:

\begin{itemize}
	\item Microsoft’s adherence to the EU-US Data Privacy Framework (DPF), adopted by the European Commission with an adequacy decision on 10 July 2023 (pursuant to Article 45 of the GDPR). The list of organisations adhering to the DPF is available on the relevant website;
	\item the conclusion of Standard Contractual Clauses (SCC) approved by the European Commission (pursuant to Article 46 of the GDPR), which Microsoft incorporates into its contractual terms as a complementary safeguard.
\end{itemize}

For further information on how Microsoft processes data, please refer to its privacy policy, available at \href{https://privacy.microsoft.com}{privacy.microsoft.com}.

\section{Rights of Data Subjects}

The data subject may exercise their rights as set out in Articles 15–22 of EU Regulation 2016/679, by contacting the Data Controller using the contact details provided in Section 1 (”Data Controller”) of this document.

Specifically, the data subject has the right, at any time, to:
\begin{itemize}
	\item \textbf{Right of access (Art. 15):} obtain confirmation of whether or not personal data concerning them exist, access such data, and receive a copy;
	\item \textbf{Right to rectification (Art. 16):} request the updating, correction of inaccurate data, or completion of incomplete data;
	\item \textbf{Right to erasure or the ”right to be forgotten” (Art. 17):} obtain the deletion of their data where they are no longer necessary for the purposes for which they were collected, or in the event of withdrawal of consent;
	\item \textbf{Right to restriction of processing (Art. 18):} request that data be solely stored (and not otherwise processed) in certain cases provided for by law;
	\item \textbf{Right to data portability (Art. 20):} receive personal data provided in a structured, commonly used, and machine-readable format, and, where technically feasible, transmit them to another controller without hindrance;
	\item \textbf{Right to object (Art. 21):} object, in whole or in part, on grounds relating to their particular situation, to the processing of personal data concerning them.
\end{itemize}

It is also noted that the selection process is not subject to decisions based solely on automated processing, including profiling, which produces legal effects or similarly significantly affects the person (Article 22 of the GDPR). Decisions regarding applications are always made by natural persons.

Without prejudice to any other administrative or judicial remedy, if the applicant considers that the processing of their data violates the provisions of the GDPR, they have the right to lodge a complaint with the Italian Data Protection Authority (Garante per la protezione dei dati personali) pursuant to Article 77 of the GDPR (\href{https://www.garanteprivacy.it}{www.garanteprivacy.it}).

Finally, in accordance with Article 7(3) and Article 6(1)(a) of the GDPR, the data subject has the right to withdraw their consent at any time, without affecting the lawfulness of processing based on consent provided prior to its withdrawal.

\end{document}
